@@
\begin{proof}
Баталгаанд энэ өгүүлбэрийн дараах леммийн (5), (6) тэнцлүүдийг ашиглана.

$a.$ $\mu$ төгсгөлөг тул компакт олонлог бүрийн хувьд $\mu(K)\leqslant\mu(\rn)<+\infty$.

$b.$ Энд $\alg=\brn$ тул энэ нөхцөл нь леммийн (5) тэнцэл юм.

Одоо $A\in\brn$, $\varepsilon>0$ байг. Леммийн (6) тэнцлээр $\mu(A\setminus C)<\dfrac{\varepsilon}{2}$ байх битүү $C\subset A$ олонлог олдоно. $K_m=C\cap[-m;m]^n$ гэвэл $K_m$ нь битүү, зааглагдсан тул компакт олонлог бөгөөд $K_1\subset K_2\subset\ldots$, $\bigcup\limits_{m=1}^{\infty}K_m=C$ байна. Хэмжээний доороос тасралтгүй чанараар $\mu(K_m)\to\mu(C)$ тул $\mu(C\setminus K_m)<\dfrac{\varepsilon}{2}$ байх $m$ олдоно. Тэгвэл $K_m\subset A$ компакт ба
$$\mu(A)-\mu(K_m)=\mu(A\setminus C)+\mu(C\setminus K_m)<\varepsilon$$
тул $\mu(A)=\sup\{\mu(K): K \ \mbox{компакт ба} \ K\subset A\}$ батлагдав.

$c.$ Энэ нөхцөл нь саяын баталсан тэнцлийн $A=U$ задгай үеийн тухайн тохиолдол юм. Иймд $\mu$ регуляр хэмжээ байна.
\end{proof}
